\documentclass[stu]{apa7}

\usepackage[spanish]{babel}
\usepackage[utf8]{inputenc}
\usepackage{hyperref}
\usepackage{enumitem}
\usepackage{geometry}
\usepackage{times}

\geometry{margin=1in}

\title{Sistema de información web contable sobre recibos de retención y movimientos de caja de la empresa minera MINCH SRL.}
\author{...}
\affiliation{...}
\course{...}
\professor{...}
\duedate{...}

\setlength{\parindent}{0.5in}

\begin{document}

\maketitle

\section{Introducción}

La informática ha evolucionado de manera acelerada desde sus inicios en la década de 1940. Las primeras computadoras eran máquinas enormes, costosas y con capacidades limitadas, utilizadas exclusivamente para fines militares y científicos. Con la invención del transistor en 1947 y posteriormente del circuito integrado en 1958, los dispositivos electrónicos comenzaron a reducir su tamaño y costo, haciendo posible su utilizacion en entornos comerciales.

En la década de 1970, el desarrollo del microprocesador fue desicivo y fundamental, permitiendo la creación de las primeras computadoras personales. Empresas como Apple e IBM popularizaron estos dispositivos, llevando la informática a hogares y pequeñas empresas. Para la década de 1990, la llegada de Internet revolucionó por completo la situacion, conectando computadoras de todo el mundo y dando origen a la era de la información.

En el siglo XXI, continuo la evolucion con la aparición de la computación en la nube, el big data, la inteligencia artificial y el Internet de las cosas. Estas tecnologías han permitido procesar volúmenes masivos de información en tiempo real, optimizando procesos y generando nuevos modelos de negocio basados en datos.

Emplear sistemas de información en el ámbito empresarial transformó radicalmente la forma de operar de las organizaciones. En sus inicios, las empresas utilizaban la informática únicamente para automatizar tareas repetitivas como la nómina, la facturación y el control de inventarios. Estos sistemas, conocidos como sistemas de procesamiento, permitieron reducir errores y aumentar la eficiencia operativa.

Con el paso del timpo, las empresas comenzaron a implementar sistemas más avanzados, como los sistemas de información gerencial, que proporcionaban informes y resúmenes para la toma de decisiones. Posteriormente, surgieron los sistemas de soporte a decisiones y los sistemas expertos, que ayudaban a los gerentes a resolver problemas complejos mediante análisis de datos y modelos predictivos.

En la actualidad, las empresas utilizan sistemas integrados, que unifican todas las areas de una organización en una sola plataforma. La ransformación digital se ha convertido en una prioridad estratégica para las empresas que buscan mantenerse competitivas en un mercado caotico

El sector minero es una de las actividades económicas más importantes a nivel mundial, especialmente en países como Perú, Chile, Australia y Canadá. La minería implica la extracción de minerales metálicos como oro, plata, cobre, hierro y no metálicos como caliza, arcilla, sal del subsuelo, los cuales son materia prima fundamental para la industria manufacturera, la construcción y la tecnología.

El mercado minero se caracteriza por su inestabilidad, ya que los precios de los minerales están sujetos a factores como la oferta y demanda global, las tensiones geopolíticas, los costos de extracción y las regulaciones ambientales. Las empresas mineras enfrentan desafíos constantes relacionados con la seguridad de sus operaciones, la gestión de residuos, el cumplimiento normativo y la relación con las comunidades locales.

En los últimos años, empresas minerías ha comenzado a adoptar tecnologías de la información para mejorar sus procesos, que incorporando sistemas de información para mejorar sus procesos internos, se han vuelto indispensables para gestionar inventarios, controlar la producción, emitir recibos de retención y procesar minerales de manera eficiente.

Tal es el caso de la empresa minera MINCH SRL. que comenzo su funcionamiento en 2023, con el fundador y propietario Wilver Churata Warachi, oficina central ubicada en ciudad de Potosi. Realizando actividades de exploracion, estraccion, procesamiento y comercializacion de minerales como zinc, plomo y plata, estos proceso en la mayoria son manejados y controlados mediante registros.

Una de las actividades realizadas dentro de la empresa es la emisión de recibos de retención. Este proceso se activa cuando la empresa adquiere bienes o servicios de un proveedor que no cuenta con NIT. En esta situación, la empresa asume la obligación de realizar el pago tributario correspondiente ante Impuestos Nacionales. El monto de dicho pago está relacionado con el tipo de bien o servicio adquirido. Actualmente, este proceso se realiza mediante hojas de Excel y se centraliza de la misma manera.

En el área de contabilidad, se realiza un proceso denominado movimientos de caja. Este proceso consiste en registrar todos los movimientos ejecutados en cada cuenta bancaria que maneja la empresa. Dichos movimientos se clasifican por cuenta bancaria y por tipo de movimiento, pudiendo ser ingresos o egresos. Actualmente, estos registros se realizan en hojas de Excel. Adicionalmente, se emite un recibo por cada movimiento realizado. Este proceso permite conocer el estado financiero de la empresa para cada tipo de cuenta bancaria que administra.

Estas actividades son realizadas por cada venta o exportación y se llevan a cabo bajo registros e informes para el seguimiento de los gastos e ingresos generados dentro de la empresa.

En la empresa MINCH SRL. se ha identificando dificultades en la gestion recibos de retencion y en los movimientos de caja, que se puede evidenciar tales como.

\begin{itemize}
    \item El área de contabilidad enfrenta dificultades significativas en el manejo de los registros de retenciones correspondientes a bienes y servicios adquiridos por la empresa. La ausencia de un sistema centralizado, sumada a la desorganización y dispersión de los registros, ocasiona problemas en el acceso oportuno a la información.

    \item El personal del área contable presenta dificultades para mantener registros exactos y actualizados sobre los movimientos realizados en cada caja. Este proceso demanda un tiempo considerable para cada tipo de caja, lo que resulta en un flujo de trabajo lento que, en ciertas circunstancias, impide su ejecución completa y oportuna.

    \item El personal contable destina aproximadamente el 40\% de su jornada laboral a tareas administrativas manuales (registro en Excel, emisión de recibos, conciliaciones), limitando el tiempo disponible para actividades de mayor valor agregado, como el análisis financiero y la planificación tributaria.

    \item La dependencia de hojas de Excel compartidas por correo electrónico o unidades de almacenamiento externo genera versiones contradictorias de la información, dificultando la trazabilidad de los cambios y la auditoría de los registros.
\end{itemize}

Por lo anteriormente mensionado, se puede evidenciar que el acceso y carencia de informacion, por la utilizacion de medios tradicionales, tambien el manejo de los registros se realizan manualmente, lo que dificulta el control y seguimiento de la informacion oportuna.

Ante este contexto, se plantea el siguiente problema de investigacion identificado:

\textbf{¿Cómo mejorar la eficiencia de la contabilidad sobre recibos de retencion y movimiento de caja de la empresa MINCH SRL.?}

Por estas circunstancias se tiene como objeto de estudio se centrara a los sistemas de informacion.

Se presenta a los sistemas de informacion web como el campo de accion que permitira el control de la contabilidad sobre recibo de retencion y movimientos de caja de la empresa minera MINCH SRL.

Para dar solucion a los problemas idenificados en la empresa minera MINCH SRL., se propone el siguiente objetivo general \textbf{implementacion de un sistemas informacion web contable sobre los recibos de retencion y movimientos de caja de la empresa minera MINCH SRL. Facilitando al personal el acceso rapido y oportuno de la informacion generada en los procesos, mejorando la eficiencia y productividad en la empresa.}

Para poder cumplir con el objetivo general se plantea los siguientes objetivos especificos:

\begin{enumerate}
    \item Desarrollar un módulo funcional para controlar y gestionar los recibos de retención, que genere reportes automáticos y reduciendo el tiempo de emisión de cada recibo.

    \item Implementar un módulo para el control y gestión de cajas, que centralice los movimientos de ingresos y egresos, permitiendo mejorar el proceso y aumentanto la productividad.

    \item Desarrollar un módulo que permite generar reportes Excel y PDF sobre el movimiento de caja, permitiendo acelerar la toma de desiciones financieras con informacion actualizada y oportuna
\end{enumerate}

Para direccionar el proyecto se plantean las siguientes Preguntas cientificas:

\begin{enumerate}
    \item ¿Cuál será la fundamentación teórica que respalda al sistema información web contable sobre recibo de retención y movimientos de caja de la empresa minera MINCH SRL?

    \item ¿Cuál es la situación actual de la información que se genera en los procesos realizados sobre los recibos de retención y movimientos de caja de la empresa minera MINCH SRL?

    \item ¿Cuál es la estrategia y metodología para el desarrollo del sistema información web contable sobre recibo de retención y movimientos de caja de la empresa minera MINCH SRL, que resuelva los problemas identificados?

    \item ¿Cuáles serán las pruebas de evaluación y calidad para que el sistema información web contable sobre recibo de retención y movimientos de caja de la empresa minera MINCH SRL cumpla su objetivo?
\end{enumerate}

Para proporcionar respuestas a las preguntas cientificas anteriormente establecidas, se proponen las siguientes tareas de investigacion:

\begin{enumerate}
    \item Elaboración de un marco teórico referencial que incluya fundamentos y definiciones mediante la revisión bibliográfica.

    \item Realizar un análisis y diagnóstico para conocer con exactitud el modelo de negocio manejado actualmente en la organización, aplicando técnicas de investigación como la entrevista y la observación.

    \item Modelación y análisis del sistema informacion web contable sobre recibos de retención y movimientos de caja para la empresa minera MINCH SRL. Se emplearán herramientas como UML y la metodología ágil SCRUM. Para el modelado del sistema, se hará uso de la herramienta Enterprise Architect para representar los diagramas y objetos del sistema o del modelo de negocio. Se utilizará MySQL como sistema gestor de base de datos, PHP como lenguaje de programación de alto nivel, y los frameworks Laravel.

    \item Realización de diferentes tecnicas de pruebas, como pruebas de funcioanlidad, pruebas de rendimiento y prueba de usuabilidad, y determinar cuales son relevantes para el sistema den desarrollo de la empresa minera MINCH SRL.
\end{enumerate}

Para la elaboracion el proyecto, se aplico los siguientes metodos a nivel teorico:

\textbf{El método de análisis-síntesis}, permite estudiar las necesidades de la empresa minera MINCH SRL, identificando los problemas y fundamentando el desarrollo del sistema. Para ello, se observaron las áreas de contabilidad, inventario y procesado de mineral. De esta manera, se logra un conocimiento integral del flujo de información en cada área, comprendiendo el problema y desglosándolo para dar una solución más apropiada.

\textbf{El método de inducción-deducción} permitió recolectar información sobre los recibos de retención. A través de este método, se analizaron etapas particulares del proceso, identificando que se trata de un proceso manual para la generación de información. Este análisis brindó datos relevantes que permitieron deducir y concretar las acciones necesarias para mejorar el proceso de los recibos de retención. De esta forma, se asegura un resultado óptimo y se brinda seguridad en la automatización del proceso, considerando los procesos relacionados.

\textbf{El método de modelación} permitirá realizar el modelado del sistema, presentando de manera exacta las tareas que se podrán ejecutar en el sistema. Mediante el uso de diagramas y figuras, se facilita la comprensión incluso para personas que no cuenten con conocimientos informáticos, asegurando una retroalimentación efectiva por parte del personal de la organización.

\textbf{El método sistémico} se utilizó para determinar el lenguaje de programación, el sistema gestor de base de datos y las herramientas que mejor interactúen durante la implementación del sistema de información para el control de inventario, recibos de retención y procesado de mineral en la empresa minera MINCH SRL.

Con los metodos empiricos, se aplicaran una se de tecnicas para recolectar informacion y poder interpretarla.

La entrevista se empleó para obtener información mediante una serie de preguntas semiestructuradas dirigidas al personal de la empresa, indagando sobre el desarrollo actual del proceso de los recibos de retención y movimientos de caja (Ver anexo 1).

La observación se empleó con el fin de obtener un panorama y una perspectiva más precisa de los procesos llevados a cabo en el área de contabilidad de la empresa minera MINCH SRL. De esta manera, se logró una visión más detallada sobre los procesos realizados por el personal encargada del área de contabilidad, evidenciando los registros que manejan (Ver anexo 1).

\textbf{Respecto a la justificación del presente proyecto}, se identifica la justificación social. Mediante la implementación del proyecto, se buscará el mejoramiento del área de contabilidad, beneficiando tanto al personal directamente involucrado como a las personas indirectamente relacionadas. Este beneficio se reflejará en un proceso más efectivo y eficiente para la emisión de recibos y movimientos de caja. De esta manera, el proyecto se justifica socialmente.

\textbf{Justificación económica} La implementación del sistema permitirá la reducción de costos en la elaboración de formularios y compra de recibos. Al sistematizar la información, se dispondrá de datos en tiempo real, accesibles y oportunos, evitando la inversión de tiempo adicional en tareas manuales.

Para el desarrollo del sistema se utilizarán herramientas de software libre. Para la implementación, solo será necesario pagar por el servicio de alojamiento del sistema en un servidor, lo que garantizará la disponibilidad del sistema desde cualquier dispositivo con conexión a internet. Esta modalidad facilita la implementación, ya que el personal del área de contabilidad cuenta con equipos con acceso a internet y con las especificaciones técnicas requeridas.

\textbf{Justificación técnica} Para el desarrollo del sistema, se utilizarán herramientas actuales de desarrollo y los conocimientos adquiridos durante la carrera, lo que aporta a la fundamentación teórica del proyecto. Se emplearán herramientas actuales y algoritmos relacionados con la contabilidad, el control de recibos de retención y los movimientos de caja, permitiendo una gestión eficiente. Además, se aprovechará la tecnología para migrar hacia un entorno digital.

Para garantizar un óptimo funcionamiento del sistema, es necesario especificar los requisitos mínimos que deben cumplir los equipos. Estos son: procesador de 1 GHz de un núcleo o superior, memoria RAM de 1 GB, contar con un sistema operativo y un navegador web que soporte HTML5.

La implementación del modelo teórico bien definido conlleva a la significación práctica en la investigación, reflejada en la utilización del sistema en la empresa minera MINCH SRL de Potosí. El uso de este sistema beneficiará significativamente al personal del área contable y a los clientes o proveedores, facilitando el acceso a la información de manera eficiente.

El aporte teórico de este estudio se centra en el diseño y análisis de un sistema de gestión avanzado que mejora la eficiencia en la administración de recibos de retención y movimientos de caja de la empresa minera MINCH SRL. Este enfoque innovador se propone optimizar la gestión de la información utilizando tecnologías informáticas avanzadas, lo que facilitará un seguimiento más preciso y en tiempo real del estado financiero, así como la exportación de información en formato PDF o Excel. Este modelo teórico brinda nuevas perspectivas sobre cómo la tecnología puede ser aplicada de manera efectiva en la gestión empresarial y establece un precedente importante para futuras implementaciones en otras áreas que requieran una gestión eficiente de datos y procesos.

La innovación científica consiste en desarrollar un sistema de información web automatizado que optimice la gestión del proceso de retención y movimientos de caja. Este sistema utiliza tecnologías avanzadas y metodologías científicas para mejorar la eficiencia, accesibilidad y fiabilidad en la administración de los recursos de la empresa.

\end{document}
